% ============================================================
%  Section 2.3.1 – Acteurs du système
%  Projet : BOA AI Hotline Automation
% ============================================================

% Packages requis (à placer dans le préambule si pas déjà présents)
% \usepackage[table,xcdraw]{xcolor}
% \usepackage{booktabs}
% \usepackage{array}
% \usepackage{tabularx}
% \usepackage{tcolorbox}
% \usepackage{fontawesome5}  % pour l'icône guideline (optionnel)

% ── Couleurs BOA ────────────────────────────────────────────
\definecolor{boaGreen}{HTML}{2E7D32}       % vert foncé entête
\definecolor{boaLightGreen}{HTML}{C8E6C9}  % vert clair lignes paires
\definecolor{boaGray}{HTML}{F5F5F5}        % gris très clair lignes impaires
\definecolor{guidelineBg}{HTML}{F1F8E9}    % fond bloc guideline
\definecolor{guidelineBorder}{HTML}{558B2F}% bordure bloc guideline

% ─────────────────────────────────────────────────────────────
\subsection{Acteurs du système}
% ─────────────────────────────────────────────────────────────

% ── Bloc Guideline ──────────────────────────────────────────
\begin{tcolorbox}[
  colback   = guidelineBg,
  colframe  = guidelineBorder,
  left      = 6pt,
  right     = 6pt,
  top       = 4pt,
  bottom    = 4pt,
  arc       = 2pt,
  boxrule   = 0.8pt
]
  \textbf{\faBookmark\; Guideline :}
  \textit{Identifiez tous les acteurs (humains et systèmes externes)
  qui interagissent avec votre application.}
\end{tcolorbox}

\vspace{6pt}

% ── Tableau des acteurs ──────────────────────────────────────
\begin{table}[h!]
  \centering
  \renewcommand{\arraystretch}{1.6}   % hauteur des lignes
  \setlength{\tabcolsep}{10pt}

  \begin{tabularx}{\textwidth}{
      >{\bfseries}p{4.5cm}   % colonne Acteur (gras)
      X                       % colonne Rôle (extensible)
    }

    % ── En-tête ────────────────────────────────────────────
    \rowcolor{boaGreen}
    \textcolor{white}{Acteur} & \textcolor{white}{Rôle} \\

    % ── Ligne 1 : Client BOA / Utilisateur ─────────────────
    \rowcolor{boaGray}
    Client BOA / Utilisateur
    &
    Acteur humain principal. Il interagit avec le \textbf{ChatBot BOA}
    pour poser des questions (FAQ), déclencher des tâches RPA bancaires
    (virement, déblocage de carte, activation de dotation E-Commerce
    ou touristique) après authentification via le \textbf{Portail BOA}.
    \\

    % ── Ligne 2 : Administrateur / Superviseur ─────────────
    \rowcolor{boaLightGreen}
    Administrateur / Superviseur
    &
    Acteur humain interne à BOA. Il accède au système via
    le \textbf{Portail Admin} pour consulter le tableau de bord,
    gérer les historiques des interactions (consulter / supprimer)
    et administrer les comptes clients BOA
    (consulter / ajouter / supprimer).
    \\

    % ── Ligne 3 : Robot RPA (UiPath) ───────────────────────
    \rowcolor{boaGray}
    Robot RPA (UiPath)
    &
    Acteur système externe. Il s'agit du robot UiPath (Robot ROA)
    chargé d'exécuter automatiquement les tâches bancaires déclenchées
    par le ChatBot : virements, déblocages de carte, activations
    de dotations. Il reçoit les instructions via le service
    d'authentification et les API REST du backend.
    \\

  \end{tabularx}

  \caption{Acteurs du système}
  \label{tab:acteurs}
\end{table}
